\documentclass[a4paper,12pt]{article}

% แพ็กเกจสำหรับภาษาไทยและการจัดการเอกสาร
\usepackage{fontspec}
\usepackage{xunicode}
\usepackage{xltxtra}
\usepackage[margin=1in]{geometry}
\usepackage{enumitem}
\usepackage{listings}
\usepackage[table]{xcolor}
\usepackage{graphicx}
\usepackage{booktabs}
\usepackage{float}
\usepackage{needspace}
\usepackage{multicol}

% ตั้งค่าฟอนต์ภาษาไทย (ต้องมี TH Sarabun New)
\setmainfont{TH Sarabun New} 
\setmonofont{Courier New}

% ปรับแต่งค่าความห่างบรรทัด
\linespread{1.15}

% ตั้งค่าสีสำหรับ Code
\definecolor{codegreen}{rgb}{0,0.6,0}
\definecolor{codegray}{rgb}{0.5,0.5,0.5}
\definecolor{codepurple}{rgb}{0.58,0,0.82}
\definecolor{backcolour}{rgb}{0.95,0.95,0.92}

\lstdefinestyle{mystyle}{
    backgroundcolor=\color{backcolour},   
    commentstyle=\color{codegreen},
    keywordstyle=\color{blue},
    numberstyle=\tiny\color{codegray},
    stringstyle=\color{codepurple},
    basicstyle=\ttfamily\footnotesize,
    breakatwhitespace=false,         
    breaklines=true,                 
    captionpos=b,                    
    keepspaces=true,                 
    showspaces=false,                
    showstringspaces=false,
    showtabs=false,                  
    tabsize=4,
    frame=single
}
\lstset{style=mystyle}

\begin{document}

% --- Cover Page / Header ---
\thispagestyle{empty}
\begin{center}
    \textbf{\Large [CPE223] COMPUTER ARCHITECTURES}\\
    \vspace{0.2cm}
    \textbf{\large M1 Mock Exam}\\
    \vspace{0.2cm}
    Date: 17 February 2026 | Time: 09:00 - 12:00 \\
    \vspace{0.5cm}
\end{center}
\hrule
\vspace{0.5cm}

% --- Part 1 ---
\section*{Part 1: Multiple Choice Questions (20 Points)}
\textbf{Instructions:} Select the correct answer.

\begin{enumerate}[label=\textbf{\arabic*.}]
    \needspace{4\baselineskip}
    \item Which statement is \textbf{true}?
    \begin{enumerate}[label=\alph*)]
        \item The motherboard contains the circuitry that connects the processor to the other hardware
        \item The term \textit{bit} refers to the amount of memory required to store one letter
        \item When you press a keyboard key, the character is stored on the computer's hard disk until it is needed by the program.
        \item If a computer has a 64-bit data bus, it can transfer 4 bytes of data at a time.
    \end{enumerate}

    \needspace{4\baselineskip}
    \item Which character is \textbf{not} Condition Codes ?
    \begin{enumerate}[label=\alph*)]
        \item N (negative)
        \item V (overflow)
        \item E (equal)
        \item Z (zero)
    \end{enumerate}

    \needspace{4\baselineskip}
    \item Which addressing mode does \textbf{not} access the main memory ?
    \begin{enumerate}[label=\alph*)]
        \item Auto-increment
        \item Relative
        \item Direct
        \item Register
    \end{enumerate}

    \needspace{4\baselineskip}
    \item The C-language statement, A = *B+1, can be compiled into (OPERATION SOURCE DES)
    \begin{enumerate}[label=\alph*)]
        \item Move 1(B), A
        \item Move (B), A; Add \#1, A
        \item Move B, A; Add \#1, A
        \item Add 1, B; Move B, A
    \end{enumerate}

    \needspace{4\baselineskip}
    \item Which statement is \textbf{false} ?
    \begin{enumerate}[label=\alph*)]
        \item 5 \hspace{0.5cm} = $(0000101)_{1'compliment}, (0000101)_{2'compliment}$
        \item -19 \hspace{0.2cm} = $(1101100)_{1'compliment}, (1101101)_{2'compliment}$
        \item -43 \hspace{0.2cm} = $(1010101)_{1'compliment}, (1010110)_{2'compliment}$
        \item -10 \hspace{0.2cm} = $(1110101)_{1'compliment}, (1110110)_{2'compliment}$
    \end{enumerate}

    \needspace{4\baselineskip}
    \item Which of the following addition generates overflow ?
    \begin{enumerate}[label=\alph*)]
        \item 11011 + 00111
        \item 10110 + 10011
        \item 11101 + 11000
        \item 00101 + 01010
    \end{enumerate}

    \needspace{4\baselineskip}
    \item Which statement is \textbf{true} ?
    \begin{enumerate}[label=\alph*)]
        \item Using a stack to save the return address of a subroutine supports subroutine nesting but not sufficient for recursion.
        \item ORIGIN 2000 \\ DATAWORD 1000 \\ The above statements cause 1000 to be placed at location 2000 during the program execution
        \item It is possible to tell whether a binary pattern represents a machine instruction or an operand.
        \item None of the above
    \end{enumerate}

    \needspace{4\baselineskip}
    \item Which statement regarding Byte Ordering (Endianness) is \textbf{TRUE}?
    \begin{enumerate}[label=\alph*)]
        \item Big-Endian stores the least significant byte (LSB) at the lowest memory address.
        \item Little-Endian stores the least significant byte (LSB) at the lowest memory address.
        \item The value 0x1234 would be stored as 34 12 in Big-Endian.
        \item Network protocols (TCP/IP) typically standardise on Little-Endian.
    \end{enumerate}

    \needspace{4\baselineskip}
    \item What is the process of removing layers from the surface of a silicon wafer during chip manufacturing called?
    \begin{enumerate}[label=\alph*)]
        \item Etching
        \item Crafting
        \item Baking
        \item Doping
    \end{enumerate}

    \needspace{4\baselineskip}
    \item Which register holds the address of the \textbf{next} instruction to be fetched?
    \begin{enumerate}[label=\alph*)]
        \item IR
        \item PC
        \item MAR
        \item MDR
    \end{enumerate}

    \needspace{4\baselineskip}
    \item Where is the instruction that is \textbf{currently being executed} stored?
    \begin{enumerate}[label=\alph*)]
        \item PC
        \item MAR
        \item IR
        \item ALU
    \end{enumerate}
    
    \needspace{4\baselineskip}
    \item Writing data to a \textbf{Special Purpose Register} (e.g., Program Counter) incorrectly is likely to cause:
    \begin{enumerate}[label=\alph*)]
        \item The computer to explode.
        \item An automatic shutdown.
        \item The hard disk to format itself.
        \item An unexpected instruction execution or branching error.
    \end{enumerate}

    \needspace{4\baselineskip}
    \item Which addressing mode is most suitable for implementing \textbf{Pointers} in high-level languages?
    \begin{enumerate}[label=\alph*)]
        \item Immediate Addressing
        \item Direct Addressing
        \item Indirect Addressing
        \item Relative Addressing
    \end{enumerate}

    \needspace{4\baselineskip}
    \item In the instruction Load R1, 16(R2), if R2 = 1000, what is the \textbf{Effective Address}?
    \begin{enumerate}[label=\alph*)]
        \item 1000
        \item 1016
        \item 16
        \item 2000
    \end{enumerate}

    \needspace{4\baselineskip}
    \item Which statement regarding Program Translation is \textbf{TRUE}?
    \begin{enumerate}[label=\alph*)]
        \item A Compiler translates assembly language into machine code.
        \item An Assembler translates assembly language into machine code.
        \item A Linker combines source code files into a text file.
        \item An Interpreter executes the high-level code directly without any translation ever happening.
    \end{enumerate}

    \needspace{4\baselineskip}
    \item What does the directive ORIGIN 200 mean?
    \begin{enumerate}[label=\alph*)]
        \item The program execution starts at line 200.
        \item It defines a constant variable named ORIGIN with value 200.
        \item The next block of code/data will be placed at memory address 200.
        \item It reserves 200 bytes of memory.
    \end{enumerate}

    \needspace{4\baselineskip}
    \item In a polling I/O mechanism, what is the function of the status flag (e.g., SIN/SOUT)?
    \begin{enumerate}[label=\alph*)]
        \item To interrupt the processor immediately.
        \item To indicate that the I/O device is ready for data transfer.
        \item To clear the input buffer.
        \item To signal that the power is on.
    \end{enumerate}

    \needspace{4\baselineskip}
    \item Which symbol represents a \textbf{Non-blocking assignment} in Verilog?
    \begin{enumerate}[label=\alph*)]
        \item =
        \item <=
        \item ==
        \item :=
    \end{enumerate}

    \needspace{4\baselineskip}
    \item Which statement about Verilog HDL is \textbf{TRUE}?
    \begin{enumerate}[label=\alph*)]
        \item Blocking assignments (=) are executed concurrently (in parallel).
        \item Non-blocking assignments (<=) are typically used for combinational logic.
        \item A flip-flop is a sequential circuit element that changes state on a clock edge.
        \item The "initial" block runs continuously throughout the simulation.
    \end{enumerate}

    \needspace{4\baselineskip}
    \item What is the meaning of always @(posedge clk)?
    \begin{enumerate}[label=\alph*)]
        \item Execute continuously when clock is 1.
        \item Execute once when clock changes from 0 to 1.
        \item Execute whenever clock changes state.
        \item Execute when clock is 0.
    \end{enumerate}
\end{enumerate}

\newpage

% --- Part 2 ---
\section*{Part 2: Short-answer questions (20 Points)}
\textbf{Instructions:} Briefly answer or explain the following concepts. (Lecture 1-3 only, No Verilog)

\begin{enumerate}[label=\textbf{\arabic*.}]
    \needspace{3\baselineskip}
    \item What does \textbf{ALU} stand for?
    \vspace{0.8cm}
    
    \needspace{3\baselineskip}
    \item What does \textbf{MDR} stand for?
    \vspace{0.8cm}
    
    \needspace{3\baselineskip}
    \item Briefly explain the function of the \textbf{Program Counter (PC)}.
    \vspace{0.8cm}

    \needspace{3\baselineskip}
    \item Briefly explain the function of the \textbf{Instruction Register (IR)}.
    \vspace{0.8cm}

    \needspace{3\baselineskip}
    \item What does \textbf{MAR} stand for?
    \vspace{0.8cm}
    
    \needspace{3\baselineskip}
    \item What does \textbf{RISC} stand for?
    \vspace{0.8cm}

    \needspace{3\baselineskip}
    \item Explain the concept of \textbf{Direct Addressing Mode}.
    \vspace{0.8cm}

    \needspace{3\baselineskip}
    \item What is the \textbf{Stack Pointer (SP)} used for?
    \vspace{0.8cm}

    \needspace{3\baselineskip}
    \item Given the value \texttt{0x12345678}, if stored in memory at address 1000 using \textbf{Little Endian} format, write down the byte value at address 1000, 1001, 1002, and 1003.
    \vspace{0.8cm}

    \needspace{3\baselineskip}
    \item Consider the following assembly loop:
    \begin{center}
    \begin{tabular}{lll}
          & MOV & R1, \#5 \\
    LOOP: & SUB & R1, \#1 \\
          & CMP & R1, \#0 \\
          & BGE & LOOP
    \end{tabular}
    \end{center}
    Explain the operation and state the \textbf{condition} that causes the loop to \textbf{terminate}.
    \vspace{1.5cm}
\end{enumerate}

\newpage

% --- Part 3 ---
\section*{Part 3: Problem Solving (20 Points)}

\textbf{1. (10 Points) The table below shows the current content of the memory and the stack. Write the NEW content in both memory and stack after execution.}

\begin{center}
    \begin{minipage}[t]{0.45\textwidth}
        \centering
        \underline{The content of the memory} \\ \vspace{0.2cm}
        \begin{tabular}{r|w{c}{1.2cm}|}
            \multicolumn{2}{c}{} \\ \cline{2-2}
            1000 & 1002 \\ \cline{2-2}
            1002 & 1004 \\ \cline{2-2}
            1004 & 1006 \\ \cline{2-2}
            1006 & 1008 \\ \cline{2-2}
            1008 & 1010 \\ \cline{2-2}
            1010 & 1012 \\ \cline{2-2}
        \end{tabular}
    \end{minipage}
    \hfill
    \begin{minipage}[t]{0.45\textwidth}
        \centering
        \underline{The content of the stack} \\ \vspace{0.2cm}
        \begin{tabular}{r|w{c}{1.2cm}|}
            \multicolumn{2}{c}{} \\ \cline{2-2}
             & \\ \cline{2-2}
             & \\ \cline{2-2}
             & \\ \cline{2-2}
            SP & 1000 \\ \cline{2-2}
             & 1002 \\ \cline{2-2}
             & 1008 \\ \cline{2-2}
        \end{tabular}
    \end{minipage}
\end{center}

\vspace{0.5cm}
After the execution of the assembly program below, write the \textbf{new content} in both memory and stack.

\vspace{0.5cm}
\begin{minipage}{0.8\textwidth}
\begin{tabular}{ll}
Load & R0, \#1004 \\
Load & R1, \#1006 \\
Load & R2, (SP) \\
Add & SP, \#1 \\
Load & R3, (SP) \\
Add & SP, \#1 \\
Subtract & SP, \#1 \\
Store & (SP), (R0)+2 \\
Add & R2, ((R1)) \\
Add & R2, (R3)+4 \\
Store & R1, R2 \\
\end{tabular}
\end{minipage}

\vspace{0.5cm}
Note: DO NOT just give the answer. Please explain how you compute the answer so that partial points can be awarded even if the final answer is wrong.

\vspace{5cm} % Space for answer

\newpage

\textbf{2. (10 Points) Translate the following C program into Assembly Language.}

\textbf{Given C Code (Filter Positive Numbers):}
\begin{center}
\fbox{\begin{minipage}{0.9\textwidth}
int N = 9;   // Index for source array 'src'
int J = 0;   // Index for destination array 'dst'
// Source Array (mix of positive and negative)
int src[10] = \{5, -20, 35, -4, 50, -6, 72, -8, 9, -100\}; \\
int dst[10]; // Destination Array
\\
while (N >= 0) \{ \\
\hspace*{0.5cm} // If element is positive, copy it to dst \\
\hspace*{0.5cm} if (src[N] > 0) \{ \\
\hspace*{1cm} dst[J] = src[N]; \\
\hspace*{1cm} J++; \\
\hspace*{0.5cm} \} \\
\hspace*{0.5cm} N--; \\
\}
\end{minipage}}
\end{center}

\textbf{Instructions:}
\begin{itemize}
    \item Use register \textbf{R1} for \texttt{N} and \textbf{R2} for \texttt{J}.
    \item Use \textbf{R3} for base address of \texttt{src} and \textbf{R4} for base address of \texttt{dst}.
    \item You must calculate the correct byte offset for both arrays using the indices.
    \item Use conditional branches.
\end{itemize}

\vspace{0.5cm}
\textbf{Answer Space:}
\vspace{8cm}

\newpage

% --- Answer Key (Has Thai) ---
\section*{Answer Key (เฉลยละเอียด)}

\subsection*{Part 1: Multiple Choice Answers}

\renewcommand{\arraystretch}{1.2}
\begin{center}
\begin{tabular}{|c|c|l|}
\hline
\rowcolor[gray]{0.8} \textbf{ข้อที่} & \textbf{เฉลย} & \textbf{คำอธิบาย (ภาษาไทย)} \\ \hline
1 & \textbf{A} & Motherboard เชื่อมต่อ Hardware ส่วนต่าง ๆ เข้าด้วยกัน \\ \hline
2 & \textbf{C} & E (Equal) เป็น condition suffix ไม่ใช่ Flag (Flags คือ N, Z, V, C) \\ \hline
3 & \textbf{D} & Register Mode อ่านข้อมูลจาก Register ใน CPU ไม่ต้องออกไป Main Memory \\ \hline
4 & \textbf{B} & Move (B), A (ย้ายค่าที่ B ชี้มา A) แล้ว Add \#1, A (บวก 1) \\ \hline
5 & \textbf{B} & ข้อ B ผิด เพราะ 1's complement ของ -19 คือ 11101100 (ไม่ใช่ 1101100) \\ \hline
6 & \textbf{B} & 10110(-10) + 10011(-13) ได้ -23 ซึ่งเกิน range 5-bit (-16) เกิด Overflow \\ \hline
7 & \textbf{B} & ORIGIN 2000 กำหนดให้ข้อมูลถัดไป (1000) เริ่มที่ Address 2000 \\ \hline
8 & \textbf{B} & Little-Endian เก็บ Byte นัยสำคัญต่ำสุด (LSB) ไว้ที่ Address ต่ำสุด \\ \hline
9 & \textbf{A} & Etching (การกัดกรด) เป็นขั้นตอนเอา Layer ออกจาก Wafer \\ \hline
10 & \textbf{B} & PC (Program Counter) เก็บ Address ของคำสั่งถัดไป \\ \hline
11 & \textbf{C} & IR (Instruction Register) เก็บคำสั่งที่กำลังถูก execute ปัจจุบัน \\ \hline
12 & \textbf{D} & การเขียนทับ Special Register อาจทำให้การทำงานผิดพลาด (Execution error) \\ \hline
13 & \textbf{C} & Indirect Addressing ใช้ค่าใน Reg/Mem เป็น Pointer ชี้ไปยังข้อมูลจริง \\ \hline
14 & \textbf{B} & Effective Address = 1000 (R2) + 16 = 1016 \\ \hline
15 & \textbf{B} & Assembler แปลง Assembly เป็น Machine Code (Compiler แปลง High-level) \\ \hline
16 & \textbf{C} & ORIGIN 200 บอกว่า Block ถัดไปให้เริ่มที่ Address 200 \\ \hline
17 & \textbf{B} & Status Flag ใช้บอกว่าอุปกรณ์ I/O พร้อมรับ/ส่งข้อมูลแล้ว (Ready) \\ \hline
18 & \textbf{B} & <= เป็นสัญลักษณ์ของ Non-blocking assignment ใน Verilog \\ \hline
19 & \textbf{C} & Flip-flop เป็นวงจร Sequential เปลี่ยนสถานะตามขอบขา Clock \\ \hline
20 & \textbf{B} & posedge clk คือช่วงที่สัญญาณนาฬิกาเปลี่ยนจาก 0 ขึ้นไป 1 \\ \hline
\end{tabular}
\end{center}

\subsection*{Part 2: Short Answer Solutions}
\begin{enumerate}
    \item \textbf{ALU}: หน่วยคำนวณทางคณิตศาสตร์และตรรกะ (Arithmetic and Logic Unit)
    \item \textbf{MDR}: Memory Data Register เก็บข้อมูลที่เพิ่งอ่านมาจากหรือกำลังจะเขียนลง Memory
    \item \textbf{PC}: เก็บ Address ของคำสั่งถัดไปที่จะถูก Fetch มาทำงาน
    \item \textbf{IR}: เก็บ Instruction ปัจจุบันที่กำลังถูก Decode/Execute อยู่
    \item \textbf{MAR}: Memory Address Register เก็บ Address ที่ต้องการติดต่อกับ Memory
    \item \textbf{RISC}: สถาปัตยกรรมคอมพิวเตอร์ที่เน้นชุดคำสั่งขนาดเล็กและเรียบง่าย (Reduced Instruction Set Computer)
    \item \textbf{Direct Addressing}: การระบุ Memory Address ลงไปตรง ๆ ในตัวคำสั่งเลย (เช่น LOAD R1, 1000)
    \item \textbf{SP}: ชี้ตำแหน่งบนสุดของ Stack (Top of Stack) ใช้เก็บข้อมูลชั่วคราวหรือ Return Address
    \item \textbf{Little Endian}: ที่ Address 1000 เก็บ Byte นัยสำคัญต่ำสุด (LSB)
    \\ Address 1000: \textbf{78}
    \\ Address 1001: \textbf{56}
    \\ Address 1002: \textbf{34}
    \\ Address 1003: \textbf{12}
    \item \textbf{Code Explanation}: ลูปนี้จะลดค่า R1 ทีละ 1 และวนซ้ำตราบเท่าที่ R1 $\ge$ 0 (วนลูปตามจำนวนรอบที่กำหนดใน R1 หรือจนกว่า R1 จะติดลบ)
\end{enumerate}

\subsection*{Part 3: Problem Solving Solutions}

\textbf{1. Memory \& Stack Content Step-by-Step Trace}
\\ \textit{สมมติฐาน: การบวก/ลบ SP ด้วย \#1 หมายถึงการเลื่อนไป Address ถัดไปตามขนาด Word (2 bytes) ในระบบจำลองนี้}

\begin{enumerate}[label=\arabic*.]
    \item \textbf{Initial State}: SP = 1000, Mem[1000]=1002, Mem[1002]=1004, ...
    \item Load R0, \#1004: R0 = 1004
    \item Load R1, \#1006: R1 = 1006
    \item Load R2, (SP): อ่าน Mem[1000] $\to$ R2 = 1002
    \item Add SP, \#1: SP เลื่อนลง $\to$ SP = 1002 (Address ถัดไป)
    \item Load R3, (SP): อ่าน Mem[1002] $\to$ R3 = 1004
    \item Add SP, \#1: SP เลื่อนลง $\to$ SP = 1004
    \item Subtract SP, \#1: SP ย้อนกลับ $\to$ SP = 1002
    \item Store (SP), (R0)+2: 
    \\ - ปลายทาง: (SP) คือ Mem[1002]
    \\ - ต้นทาง: ใช้ R0 (1004) เป็น Address $\to$ อ่าน Mem[1004] ได้ค่า \textbf{1006}
    \\ - หลังอ่าน R0 บวอกเพิ่ม 2 $\to$ R0 = 1006 (Post-increment)
    \\ - \textbf{ผลลัพธ์}: เขียนค่า 1006 ทับลงใน Mem[1002] (ค่าเดิมคือ 1004 หายไป)
    \item Add R2, ((R1)):
    \\ - R1 = 1006
    \\ - Value = Mem[Mem[1006]] = Mem[1008] = 1010
    \\ - R2 = R2 + 1010 = 1002 + 1010 = \textbf{2012}
    \item Add R2, (R3)+4:
    \\ - R3 ปัจจุบันคือ 1004 (โหลดไว้ตั้งแต่ข้อ 6 ก่อน Mem[1002] จะเปลี่ยน และ R3 ไม่ได้ถูกแก้)
    \\ - อ่าน Mem[R3] = Mem[1004] = 1006
    \\ - R2 = R2 + 1006 = 2012 + 1006 = \textbf{3018}
    \\ - R3 บวกเพิ่ม 4 $\to$ R3 = 1008
    \item Store R1, R2:
    \\ - เขียนค่า R2 (3018) ลงใน Address ที่ R1 ชี้ (1006)
    \\ - \textbf{ผลลัพธ์}: Mem[1006] เปลี่ยนเป็น 3018
\end{enumerate}

\textbf{Final Answers:}
\begin{center}
    \begin{minipage}[t]{0.3\textwidth}
        \centering
        \underline{Memory} \\
        \begin{tabular}{|c|c|}
            \hline
            1000 & 1002 \\ \hline
            1002 & \textbf{1006} \\ \hline
            1004 & 1006 \\ \hline
            1006 & \textbf{3018} \\ \hline
            1008 & 1010 \\ \hline
            1010 & 1012 \\ \hline
        \end{tabular}
    \end{minipage}
    \hspace{1cm}
    \begin{minipage}[t]{0.3\textwidth}
        \centering
        \underline{Stack} \\
        (SP อยู่ที่ 1002) \\
        \begin{tabular}{|c|}
            \hline
             \\ \hline
             \\ \hline
            1000 \\ \hline
            \textbf{1002} \\ \hline
            1008 \\ \hline
        \end{tabular}
    \end{minipage}
\end{center}

\vspace{0.5cm}
\textbf{2. C to Assembly Translation (Filter Positive Numbers)}
\\ \textbf{Logic }: ใช้ R1(N), R2(J), R3(Base src), R4(Base dst). ต้องคำนวณ Offset แยกกัน 2 ตัว.

\begin{center}
\begin{tabular}{llll}
SRC\_DATA & DCD   & 5, -20, 35, -4, 50, -6, 72, -8, 9, -100 & ; Source Array \\
DST\_DATA & SPACE & 40                 & ; Allocate 40 bytes (10 ints) \\
\\
START:    & LDR   & R1, =9            & ; N = 9 (Index for src) \\
          & MOV   & R2, \#0            & ; J = 0 (Index for dst) \\
          & LDR   & R3, =SRC\_DATA     & ; Base Address src \\
          & LDR   & R4, =DST\_DATA     & ; Base Address dst \\
\\
LOOP:     & CMP   & R1, \#0            & ; Loop Condition: N >= 0 \\
          & BLT   & DONE              & ; If N < 0, Exit Loop \\
\\
          & LSL   & R5, R1, \#2        & ; Calculate Offset for src: R5 = N * 4 \\
          & LDR   & R6, [R3, R5]      & ; Load src[N] into R6 \\
\\
          & CMP   & R6, \#0            & ; Check if src[N] > 0 \\
          & BLE   & SKIP\_COPY         & ; If negative or zero, skip copy \\
\\
          & LSL   & R7, R2, \#2        & ; Calculate Offset for dst: R7 = J * 4 \\
          & STR   & R6, [R4, R7]      & ; Store src[N] into dst[J] \\
          & ADD   & R2, R2, \#1        & ; Increment J (J++) \\
\\
SKIP\_COPY: & SUB   & R1, R1, \#1        & ; Decrement N (N--) \\
          & B     & LOOP              & ; Repeat Loop \\
\\
DONE:     & ; Program End             & \\
\end{tabular}
\end{center}
\end{document}